% GENERATED FILE. Edit canonical lessons, then run npm run generate:books.
% canonical-source-hash: fnv1a64:0e976893717b4b68
% canonical-lessons: JA-C412-touge, JA-C412-choujou, JA-C412-sanmyaku, JA-C412-kougen, JA-C412-tozan

\chapter{Mountain Passes and High Plains}
\label{ch:ja-mountain-passes-and-high-plains}
\hlchaptermodality{412}

\begin{chapteropening}
\textbf{By the end of this chapter:} \emph{I can talk about climbing in the mountains.}

Complete the last lesson of chapter 412: I can talk about climbing in the mountains.
\end{chapteropening}

\section[\texorpdfstring{\ja{とうげ} (tōge)}{とうげ (tōge)}]{\ja{とうげ} (tōge) — a mountain pass}

\label{lesson:JA-C412-touge}



\begin{quote}
Before the new one: say the Japanese for to surround, then the Japanese for to scoop up.
\end{quote}

\subsection*{You'll want to know: \ja{とうげ}}
\textbf{\ja{とうげ}} — \emph{tōge} — \textquotedblleft{}a mountain pass\textquotedblright{}.

\begin{quote}
\textbf{\ja{くるまで} \ja{とうげを} \ja{こえました。}} — \emph{kuruma de tōge o koemashita} — \textquotedblleft{}We crossed the mountain pass by car.\textquotedblright{}
\end{quote}

\begin{cousinweb}
\textbf{\ja{とうげ}} (\emph{tōge}) is a native Japanese word (wago), inherited rather than borrowed from Chinese; its kanji was attached later as a kun'yomi (JMdict).
\end{cousinweb}

\begin{grammarlens}[title={What you've built}]
One more word for getting around.
\end{grammarlens}

\subsection*{Your turn}
\begin{itemize}
\raggedright
  \item \emph{Say it:} \emph{tōge}
  \item \emph{Say it:} \emph{tōge}, once more
  \item \emph{Say it:} \emph{kuruma de tōge o koemashita} — We crossed the mountain pass by car.
  \item \emph{Recall:} say the Japanese for to surround, then the Japanese for to scoop up, then say \emph{tōge} again
\end{itemize}

\begin{tcolorbox}[breakable,colback=teal!4,colframe=teal!35!black,title={Before you move on}]
What does \ja{とうげ} mean? (\textquotedblleft{}A mountain pass\textquotedblright{}.) Say it once more.
\end{tcolorbox}
\section[\texorpdfstring{\ja{ちょうじょう} (chōjō)}{ちょうじょう (chōjō)}]{\ja{ちょうじょう} (chōjō) — a summit; the top}

\label{lesson:JA-C412-choujou}



\begin{quote}
Before the new one: say the Japanese for to scoop up, then the Japanese for a mountain pass.
\end{quote}

\subsection*{You'll want to know: \ja{ちょうじょう}}
\textbf{\ja{ちょうじょう}} — \emph{chōjō} — \textquotedblleft{}a summit; the top\textquotedblright{}.

\begin{quote}
\textbf{\ja{やまの} \ja{ちょうじょうは} \ja{さむかったです。}} — \emph{yama no chōjō wa samukatta desu} — \textquotedblleft{}It was cold at the top of the mountain.\textquotedblright{}
\end{quote}

\begin{cousinweb}
\textbf{\ja{ちょうじょう}} (\emph{chōjō}) is Sino-Japanese (kango): two Chinese characters read with their on'yomi, readings borrowed from Chinese. They are \emph{chō}, \textquotedblleft{}summit\textquotedblright{} (Middle Chinese \emph{tengX}) and \emph{jō}, \textquotedblleft{}above\textquotedblright{} (\emph{dzyangH}) (KANJIDIC2; Middle Chinese in Baxter's transcription).
\end{cousinweb}

\begin{grammarlens}[title={What you've built}]
One more word for getting around.
\end{grammarlens}

\subsection*{Your turn}
\begin{itemize}
\raggedright
  \item \emph{Say it:} \emph{chōjō}
  \item \emph{Say it:} \emph{chōjō}, once more
  \item \emph{Say it:} \emph{yama no chōjō wa samukatta desu} — It was cold at the top of the mountain.
  \item \emph{Recall:} say the Japanese for to scoop up, then the Japanese for a mountain pass, then say \emph{chōjō} again
\end{itemize}

\begin{tcolorbox}[breakable,colback=teal!4,colframe=teal!35!black,title={Before you move on}]
What does \ja{ちょうじょう} mean? (\textquotedblleft{}A summit; the top\textquotedblright{}.) Say it once more.
\end{tcolorbox}
\section[\texorpdfstring{\ja{さんみゃく} (sanmyaku)}{さんみゃく (sanmyaku)}]{\ja{さんみゃく} (sanmyaku) — a mountain range}

\label{lesson:JA-C412-sanmyaku}



\begin{quote}
Before the new one: say the Japanese for a mountain pass, then the Japanese for a summit.
\end{quote}

\subsection*{You'll want to know: \ja{さんみゃく}}
\textbf{\ja{さんみゃく}} — \emph{sanmyaku} — \textquotedblleft{}a mountain range\textquotedblright{}.

\begin{quote}
\textbf{\ja{ひこうきから} \ja{さんみゃくが} \ja{みえました。}} — \emph{hikōki kara sanmyaku ga miemashita} — \textquotedblleft{}We could see the mountain range from the plane.\textquotedblright{}
\end{quote}

\begin{cousinweb}
\textbf{\ja{さんみゃく}} (\emph{sanmyaku}) is Sino-Japanese (kango): two Chinese characters read with their on'yomi, readings borrowed from Chinese. They are \emph{san}, \textquotedblleft{}mountain\textquotedblright{} (Middle Chinese \emph{srean}) and \emph{myaku}, \textquotedblleft{}vein\textquotedblright{} (\emph{meak}) (KANJIDIC2; Middle Chinese in Baxter's transcription).
\end{cousinweb}

\begin{grammarlens}[title={What you've built}]
One more word for getting around.
\end{grammarlens}

\subsection*{Your turn}
\begin{itemize}
\raggedright
  \item \emph{Say it:} \emph{sanmyaku}
  \item \emph{Say it:} \emph{sanmyaku}, once more
  \item \emph{Say it:} \emph{hikōki kara sanmyaku ga miemashita} — We could see the mountain range from the plane.
  \item \emph{Recall:} say the Japanese for a mountain pass, then the Japanese for a summit, then say \emph{sanmyaku} again
\end{itemize}

\begin{tcolorbox}[breakable,colback=teal!4,colframe=teal!35!black,title={Before you move on}]
What does \ja{さんみゃく} mean? (\textquotedblleft{}A mountain range\textquotedblright{}.) Say it once more.
\end{tcolorbox}
\section[\texorpdfstring{\ja{こうげん} (kōgen)}{こうげん (kōgen)}]{\ja{こうげん} (kōgen) — a plateau; highlands}

\label{lesson:JA-C412-kougen}



\begin{quote}
Before the new one: say the Japanese for a summit, then the Japanese for a mountain range.
\end{quote}

\subsection*{You'll want to know: \ja{こうげん}}
\textbf{\ja{こうげん}} — \emph{kōgen} — \textquotedblleft{}a plateau; highlands\textquotedblright{}.

\begin{quote}
\textbf{\ja{なつは} \ja{こうげんが} \ja{すずしいです。}} — \emph{natsu wa kōgen ga suzushii desu} — \textquotedblleft{}The highlands are cool in summer.\textquotedblright{}
\end{quote}

\begin{cousinweb}
\textbf{\ja{こうげん}} (\emph{kōgen}) is Sino-Japanese (kango): two Chinese characters read with their on'yomi, readings borrowed from Chinese. They are \emph{kō}, \textquotedblleft{}high\textquotedblright{} (Middle Chinese \emph{kaw}) and \emph{gen}, \textquotedblleft{}plain\textquotedblright{} (\emph{ngjwon}) (KANJIDIC2; Middle Chinese in Baxter's transcription).
\end{cousinweb}

\begin{grammarlens}[title={What you've built}]
One more word for getting around.
\end{grammarlens}

\subsection*{Your turn}
\begin{itemize}
\raggedright
  \item \emph{Say it:} \emph{kōgen}
  \item \emph{Say it:} \emph{kōgen}, once more
  \item \emph{Say it:} \emph{natsu wa kōgen ga suzushii desu} — The highlands are cool in summer.
  \item \emph{Recall:} say the Japanese for a summit, then the Japanese for a mountain range, then say \emph{kōgen} again
\end{itemize}

\begin{tcolorbox}[breakable,colback=teal!4,colframe=teal!35!black,title={Before you move on}]
What does \ja{こうげん} mean? (\textquotedblleft{}A plateau; highlands\textquotedblright{}.) Say it once more.
\end{tcolorbox}
\section[\texorpdfstring{\ja{とざん} (tozan)}{とざん (tozan)}]{\ja{とざん} (tozan) — mountain climbing}

\label{lesson:JA-C412-tozan}



\begin{quote}
Before the new one: say the Japanese for a mountain range, then the Japanese for a plateau.
\end{quote}

\subsection*{You'll want to know: \ja{とざん}}
\textbf{\ja{とざん}} — \emph{tozan} — \textquotedblleft{}mountain climbing\textquotedblright{}.

\begin{quote}
\textbf{\ja{しゅうまつに} \ja{ちちと} \ja{とざんに} \ja{いきました。}} — \emph{shūmatsu ni chichi to tozan ni ikimashita} — \textquotedblleft{}At the weekend I went mountain climbing with my father.\textquotedblright{}
\end{quote}

\begin{cousinweb}
\textbf{\ja{とざん}} (\emph{tozan}) is Sino-Japanese (kango): two Chinese characters read with their on'yomi, readings borrowed from Chinese. They are \emph{to}, \textquotedblleft{}ascend\textquotedblright{} (Middle Chinese \emph{tong}) and \emph{zan}, \textquotedblleft{}mountain\textquotedblright{} (\emph{srean}) (KANJIDIC2; Middle Chinese in Baxter's transcription).
\end{cousinweb}

\begin{grammarlens}[title={What you've built}]
One more word for getting around.
\end{grammarlens}

\subsection*{Your turn}
\begin{itemize}
\raggedright
  \item \emph{Say it:} \emph{tozan}
  \item \emph{Say it:} \emph{tozan}, once more
  \item \emph{Say it:} \emph{shūmatsu ni chichi to tozan ni ikimashita} — At the weekend I went mountain climbing with my father.
  \item \emph{Recall:} say the Japanese for a mountain range, then the Japanese for a plateau, then say \emph{tozan} again
\end{itemize}

\begin{tcolorbox}[breakable,colback=teal!4,colframe=teal!35!black,title={Before you move on}]
What does \ja{とざん} mean? (\textquotedblleft{}Mountain climbing\textquotedblright{}.) Say it once more.
\end{tcolorbox}
